\chapter{Curves over a field}
\label{chap:Curves}

Throughout this chapter \(k\) is a field. As in Chapter~\ref{chap:Challenge}, a \emph{smooth
curve} over \(k\) is a \(k\)-scheme which is smooth of relative dimension one, proper, and
geometrically irreducible over \(k\). This chapter provides the geometric substrate for the
genus: smooth schemes over a field are reduced, so a smooth curve is geometrically integral;
a proper geometrically integral scheme over a field has no global functions besides the
constants; the projective line carries two affine charts with polynomial section rings and
Laurent-polynomial overlap; and, by Zariski's main theorem, producing a finite morphism
\(C \to \PP^1\) reduces to producing a locally quasi-finite one.

\section{Smooth algebras over a domain are reduced}

\begin{lemma}[Unramified flat algebras over a domain are reduced]
  \label{lem:unramified_flat_reduced}
  \lean{Algebra.FormallyUnramified.isReduced_of_isDomain}
  \leanok
  Let \(B\) be an integral domain and let \(S\) be a commutative \(B\)-algebra which is
  formally unramified, essentially of finite type, and flat over \(B\). Then \(S\) is reduced.
\end{lemma}
\begin{proof}
  \leanok
  Let \(K = \Frac B\) be the fraction field. The \(K\)-algebra \(K \otimes_B S\) is again
  formally unramified and essentially of finite type, both properties being stable under base
  change, and a formally unramified, essentially of finite type algebra over a field is
  reduced. It therefore suffices to embed \(S\) into \(K \otimes_B S\). The canonical map
  \(s \mapsto 1 \otimes s\) is, up to the isomorphism \(S \cong B \otimes_B S\), obtained from
  the injection \(B \hookrightarrow K\) by tensoring with \(S\); since \(S\) is flat over
  \(B\), it is injective. A subring of a reduced ring is reduced.
\end{proof}

\begin{lemma}[\'Etale algebras over a domain are reduced]
  \label{lem:etale_reduced_domain}
  \lean{Algebra.Etale.isReduced_of_isDomain}
  \leanok
  An \'etale algebra over an integral domain is reduced.
\end{lemma}
\begin{proof}
  \leanok
  \uses{lem:unramified_flat_reduced}
  An \'etale algebra is formally unramified and of finite presentation, in particular
  essentially of finite type, and it is flat. Apply \ref{lem:unramified_flat_reduced}.
\end{proof}

\begin{lemma}[Standard smooth homomorphisms out of a domain]
  \label{lem:standard_smooth_reduced_ringhom}
  \lean{RingHom.IsStandardSmooth.isReduced_of_isDomain}
  \leanok
  Let \(B\) be an integral domain and let \(B \to S\) be a standard smooth ring homomorphism,
  that is, one admitting a presentation \(S = B[x_1,\dots,x_n]/(f_1,\dots,f_c)\) in which the
  Jacobian matrix \((\partial f_u/\partial x_v)_{u,v \le c}\) has invertible determinant in
  \(S\). Then \(S\) is reduced.
\end{lemma}
\begin{proof}
  \leanok
  \uses{lem:etale_reduced_domain}
  By the structure theorem for standard smooth algebras, \(S\) is \'etale over a polynomial
  ring over the base: there are \(m\) and a factorization
  \(B \to B[x_1,\dots,x_m] \to S\) with the second map \'etale. A polynomial ring over an
  integral domain is an integral domain, so \ref{lem:etale_reduced_domain} applies.
\end{proof}

\begin{lemma}[Standard smooth algebras over a domain]
  \label{lem:standard_smooth_reduced}
  \lean{Algebra.IsStandardSmooth.isReduced_of_isDomain}
  \leanok
  A standard smooth algebra over an integral domain is reduced.
\end{lemma}
\begin{proof}
  \leanok
  \uses{lem:standard_smooth_reduced_ringhom}
  The structure homomorphism is a standard smooth ring homomorphism out of a domain; apply
  \ref{lem:standard_smooth_reduced_ringhom}.
\end{proof}

\begin{theorem}[Smooth schemes over a field are reduced]
  \label{thm:smooth_reduced_field}
  \lean{AlgebraicGeometry.Smooth.isReduced_of_field}
  \leanok
  Let \(K\) be a field and let \(X\) be a scheme which is smooth over \(\Spec K\). Then \(X\)
  is reduced.
\end{theorem}
\begin{proof}
  \leanok
  \uses{lem:standard_smooth_reduced_ringhom}
  A smooth morphism is, locally on the source and the target, standard smooth: every point
  \(x \in X\) has an affine open neighbourhood \(V\), mapping into an affine open \(U\) of the
  base, such that the induced ring homomorphism \(\Gamma(U) \to \Gamma(V)\) is standard
  smooth. Since \(\Spec K\) has a single point and \(U\) is a nonempty open subset containing
  the image of \(x\), we have \(U = \Spec K\), and \(\Gamma(U) \cong K\) is a field, in
  particular a domain. By \ref{lem:standard_smooth_reduced_ringhom} the ring \(\Gamma(V)\) is
  reduced, hence the affine scheme \(V\) is reduced. The opens \(V\) cover \(X\), and a scheme
  covered by reduced open subschemes is reduced.
\end{proof}

\begin{theorem}[Smooth morphisms are geometrically reduced]
  \label{thm:smooth_geometrically_reduced}
  \lean{AlgebraicGeometry.Smooth.geometricallyReduced}
  \leanok
  A smooth morphism of schemes \(f : X \to Y\) is geometrically reduced: for every field \(K\)
  and every morphism \(\Spec K \to Y\), the base change \(X \times_Y \Spec K\) is a reduced
  scheme.
\end{theorem}
\begin{proof}
  \leanok
  \uses{thm:smooth_reduced_field}
  Smoothness is stable under base change, so the projection
  \(X \times_Y \Spec K \to \Spec K\) is smooth; apply \ref{thm:smooth_reduced_field}.
\end{proof}

\begin{lemma}[Smooth of relative dimension one implies smooth]
  \label{lem:smooth_of_relDim_one}
  \lean{AlgebraicGeometry.Smooth.of_smoothOfRelativeDimension_one}
  \leanok
  A morphism of schemes which is smooth of relative dimension one is smooth.
\end{lemma}
\begin{proof}
  \leanok
  A morphism smooth of some relative dimension is smooth.
\end{proof}

\begin{lemma}[Smooth of a relative dimension implies geometrically reduced]
  \label{lem:relDim_geometrically_reduced}
  \lean{AlgebraicGeometry.SmoothOfRelativeDimension.geometricallyReduced}
  \leanok
  For every \(n\), a morphism of schemes which is smooth of relative dimension \(n\) is
  geometrically reduced.
\end{lemma}
\begin{proof}
  \leanok
  \uses{thm:smooth_geometrically_reduced}
  Such a morphism is smooth; apply \ref{thm:smooth_geometrically_reduced}.
\end{proof}

\begin{theorem}[Smooth and geometrically irreducible implies geometrically integral]
  \label{thm:smooth_geometrically_integral}
  \lean{AlgebraicGeometry.Smooth.geometricallyIntegral}
  \leanok
  A smooth, geometrically irreducible morphism of schemes is geometrically integral: every
  base change to the spectrum of a field is an integral scheme.
\end{theorem}
\begin{proof}
  \leanok
  \uses{thm:smooth_geometrically_reduced}
  Every base change to the spectrum of a field is reduced
  (\ref{thm:smooth_geometrically_reduced}) and irreducible (geometric irreducibility), and a
  reduced scheme with irreducible --- in particular nonempty --- underlying space is integral.
\end{proof}

\begin{theorem}[Relative-dimension form]
  \label{thm:relDim_geometrically_integral}
  \lean{AlgebraicGeometry.SmoothOfRelativeDimension.geometricallyIntegral}
  \leanok
  For every \(n\), a geometrically irreducible morphism of schemes which is smooth of relative
  dimension \(n\) is geometrically integral.
\end{theorem}
\begin{proof}
  \leanok
  \uses{thm:smooth_geometrically_integral}
  Such a morphism is smooth; apply \ref{thm:smooth_geometrically_integral}.
\end{proof}

\section{Schemes over a field: the topological substrate}

In this section \(X\) is a scheme over \(k\), that is, a scheme equipped with a structure
morphism \(X \to \Spec k\).

\begin{lemma}[Irreducibility]
  \label{lem:irreducible_of_geomirred}
  \lean{AlgebraicGeometry.irreducibleSpace_left_of_geometricallyIrreducible}
  \leanok
  If \(X\) is geometrically irreducible over \(k\), then the underlying topological space of
  \(X\) is irreducible; in particular \(X\) is nonempty.
\end{lemma}
\begin{proof}
  \leanok
  Geometric irreducibility applied to the trivial extension \(k \to k\) says that
  \(X \times_{\Spec k} \Spec k \cong X\) is irreducible.
\end{proof}

\begin{lemma}[Quasi-compactness]
  \label{lem:curve_compact}
  \lean{AlgebraicGeometry.compactSpace_left_of_quasiCompact}
  \leanok
  If \(X\) is quasi-compact over \(\Spec k\) (for instance proper), then the underlying
  topological space of \(X\) is quasi-compact.
\end{lemma}
\begin{proof}
  \leanok
  \(\Spec k\) is a one-point space, hence quasi-compact, and \(X\) is the preimage of it
  under a quasi-compact morphism.
\end{proof}

\begin{lemma}[Quasi-separatedness]
  \label{lem:curve_quasiseparated}
  \lean{AlgebraicGeometry.quasiSeparatedSpace_left_of_quasiSeparated}
  \leanok
  If \(X\) is quasi-separated over \(\Spec k\) (for instance proper), then the underlying
  topological space of \(X\) is quasi-separated: the intersection of any two quasi-compact
  opens is quasi-compact.
\end{lemma}
\begin{proof}
  \leanok
  \(\Spec k\) is affine, hence a quasi-separated scheme, and a scheme quasi-separated over a
  quasi-separated scheme is quasi-separated.
\end{proof}

\begin{proposition}[Integrality]
  \label{prop:curve_integral}
  \lean{AlgebraicGeometry.isIntegral_left_of_geometricallyReduced}
  \leanok
  If \(X\) is geometrically reduced and geometrically irreducible over \(k\), then \(X\) is an
  integral scheme.
\end{proposition}
\begin{proof}
  \leanok
  \uses{lem:irreducible_of_geomirred}
  Applying geometric reducedness to the trivial extension \(k \to k\) shows that \(X\) is
  reduced, and \(X\) is irreducible by \ref{lem:irreducible_of_geomirred}. A reduced scheme
  with irreducible underlying space is integral.
\end{proof}

\section{Global sections of a proper geometrically integral scheme}

\begin{lemma}[Pushout of a field extension along itself]
  \label{lem:pushout_field_epi}
  \lean{CommRingCat.epi_of_isPushout_of_isField}
  \leanok
  Let \(\varphi : A \to B\) and \(\psi : A \to C\) be homomorphisms of commutative rings with
  \(B\) nontrivial, and let \(D\), with coprojections \(\iota_B : B \to D\) and
  \(\iota_C : C \to D\), be a pushout of \((\varphi, \psi)\) (for instance
  \(D = B \otimes_A C\)). Assume that \(D\) is a field and that there is a ring isomorphism
  \(e : C \cong B\) with \(e \circ \psi = \varphi\). Then \(\varphi\) is an epimorphism in the
  category of commutative rings.
\end{lemma}
\begin{proof}
  \leanok
  Let \(\mu : D \to B\) be the homomorphism with \(\mu \circ \iota_B = \id_B\) and
  \(\mu \circ \iota_C = e\), given by the universal property of the pushout (the two maps
  agree on \(A\) because \(e \circ \psi = \varphi\)). Since \(D\) is a field and \(B\) is
  nontrivial, \(\mu\) is injective; it is surjective because \(\mu \circ \iota_B = \id_B\).
  Hence \(\mu\) is an isomorphism, and from
  \(\mu \circ \iota_B = \id_B = \mu \circ (\iota_C \circ e^{-1})\) we conclude
  \(\iota_B = \iota_C \circ e^{-1}\).

  Now let \(u, v : B \to W\) satisfy \(u \circ \varphi = v \circ \varphi\). Then
  \(u \circ \varphi = (v \circ e) \circ \psi\), so there is
  \(w : D \to W\) with \(w \circ \iota_B = u\) and \(w \circ \iota_C = v \circ e\). Therefore
  \[
    u = w \circ \iota_B = w \circ \iota_C \circ e^{-1} = v \circ e \circ e^{-1} = v,
  \]
  which proves that \(\varphi\) is an epimorphism.
\end{proof}

\begin{lemma}[Global sections form a field]
  \label{lem:sections_field}
  \lean{AlgebraicGeometry.isField_of_universallyClosed_of_isField}
  \leanok
  Let \(R\) be a commutative ring which is a field, let \(X\) be an integral scheme, and let
  \(f : X \to \Spec R\) be universally closed. Then \(\Gamma(X, \struct X)\) is a field.
\end{lemma}
\begin{proof}
  \leanok
  \(\Gamma(X, \struct X)\) is an integral domain because \(X\) is integral. Since \(f\) is
  universally closed, the induced ring homomorphism
  \(R \cong \Gamma(\Spec R, \struct{\Spec R}) \to \Gamma(X, \struct X)\) is an integral
  extension. A domain which is integral over a field is a field.
\end{proof}

\begin{theorem}[Constant functions on a proper geometrically integral scheme]
  \label{thm:sections_k}
  \lean{AlgebraicGeometry.bijective_appTop_of_isProper_of_geometricallyIntegral}
  \leanok
  Let \(K\) be a field and let \(X\) be a scheme, proper and geometrically integral over
  \(K\). Then the ring homomorphism
  \(\Gamma(\Spec K, \struct{\Spec K}) \to \Gamma(X, \struct X)\) induced by the structure
  morphism is bijective; identifying \(\Gamma(\Spec K, \struct{\Spec K}) = K\), this says
  \(\Gamma(X, \struct X) = K\).
\end{theorem}
\begin{proof}
  \leanok
  \uses{lem:pushout_field_epi, lem:sections_field, lem:curve_compact, lem:curve_quasiseparated}
  Write \(A = \Gamma(X, \struct X)\) and let \(\varphi : K \to A\) be the map in question.
  Applying geometric integrality to the trivial extension \(K \to K\) shows that \(X\) is
  integral; in particular \(X\) is nonempty and \(A\) is a nontrivial domain, so \(\varphi\),
  a ring homomorphism out of a field into a nontrivial ring, is injective. By
  \ref{lem:sections_field} (properness includes universal closedness), \(A\) is a field, and
  since \(f\) is proper, \(A\) is moreover a finite \(K\)-module. As a finite ring
  homomorphism is surjective as soon as it is an epimorphism of commutative rings, it remains
  to prove that \(\varphi\) is an epimorphism.

  Let \(q : \Spec A \to \Spec K\) be the morphism induced by \(\varphi\) and form the base
  change \(Y = X \times_{\Spec K} \Spec A\). Since \(A\) is a field, geometric integrality of
  \(f\) makes \(Y\) integral, and the projection \(Y \to \Spec A\) is universally closed as a
  base change of \(f\); hence \(D := \Gamma(Y, \struct Y)\) is a field by
  \ref{lem:sections_field}. On the other hand, \(X\) is quasi-compact and quasi-separated as
  a topological space (\ref{lem:curve_compact}, \ref{lem:curve_quasiseparated}), and \(q\) is
  flat because every vector space over \(K\) is flat; flat base change for global sections
  therefore exhibits the square of section rings
  \[
    \begin{array}{ccc}
      K & \longrightarrow & A \\
      \downarrow & & \downarrow \\
      A & \longrightarrow & D
    \end{array}
  \]
  (left column \(\varphi\), top row \(\varphi\) transported through the canonical isomorphism
  \(\Gamma(\Spec A, \struct{\Spec A}) \cong A\)) as a pushout square of commutative rings,
  \(D \cong A \otimes_K A\). Its pushout \(D\) is a field and its two legs are isomorphic
  copies of \(\varphi\), so \(\varphi\) is an epimorphism by \ref{lem:pushout_field_epi}.
\end{proof}

\begin{corollary}[Isomorphism form]
  \label{cor:sections_iso}
  \lean{AlgebraicGeometry.isIso_appTop_of_isProper_of_geometricallyIntegral}
  \leanok
  In the situation of \ref{thm:sections_k}, the induced map
  \(\Gamma(\Spec K, \struct{\Spec K}) \to \Gamma(X, \struct X)\) is an isomorphism of rings.
\end{corollary}
\begin{proof}
  \leanok
  \uses{thm:sections_k}
  A bijective ring homomorphism is an isomorphism.
\end{proof}

\begin{corollary}[Global sections of the curve]
  \label{cor:curve_sections}
  \lean{AlgebraicGeometry.isIso_hom_appTop_of_geometricallyReduced}
  \leanok
  Let \(C\) be a scheme over \(k\) which is proper, geometrically irreducible and
  geometrically reduced over \(k\). Then the structure morphism induces an isomorphism
  \(\Gamma(\Spec k, \struct{\Spec k}) \cong \Gamma(C, \struct C)\); that is,
  \(\Gamma(C, \struct C) = k\).
\end{corollary}
\begin{proof}
  \leanok
  \uses{cor:sections_iso}
  Every base change of \(C\) to the spectrum of a field is reduced and irreducible, hence
  integral; so \(C\) is geometrically integral over \(k\) and \ref{cor:sections_iso} applies.
\end{proof}

\begin{remark}[The smooth curve]
  \label{rem:curve_bundle}
  \uses{lem:smooth_of_relDim_one, lem:relDim_geometrically_reduced,
    thm:relDim_geometrically_integral, lem:irreducible_of_geomirred, lem:curve_compact,
    lem:curve_quasiseparated, prop:curve_integral, cor:curve_sections}
  \leanok
  Let \(C\) be a smooth curve over \(k\), i.e.\ proper, smooth of relative dimension one and
  geometrically irreducible over \(k\), as in Chapter~\ref{chap:Challenge}. Combining the
  above: \(C\) is smooth over \(k\) (\ref{lem:smooth_of_relDim_one}), geometrically reduced
  (\ref{lem:relDim_geometrically_reduced}) and geometrically integral
  (\ref{thm:relDim_geometrically_integral}); its underlying space is irreducible
  (\ref{lem:irreducible_of_geomirred}), quasi-compact (\ref{lem:curve_compact}) and
  quasi-separated (\ref{lem:curve_quasiseparated}); \(C\) is an integral scheme
  (\ref{prop:curve_integral}); and \(\Gamma(C, \struct C) = k\) (\ref{cor:curve_sections}).
\end{remark}

\section{The projective line and its charts}

Let \(A = k[X_0, X_1]\) be the polynomial ring in two variables, graded by total degree. For a
homogeneous element \(f \in A\) we write \(A_{(f)}\) for the degree-zero part of the
localization \(A_f\); its elements are fractions \(a/f^n\) with \(a\) homogeneous of degree
\(n \deg f\).

\begin{definition}[The projective line]
  \label{def:P1}
  \lean{AlgebraicGeometry.P1}
  \leanok
  The projective line over \(k\) is \(\PP^1 = \Proj A\), the Proj of the graded ring
  \(A = k[X_0, X_1]\).
\end{definition}

\begin{definition}[Degree zero]
  \label{def:P1_gradeZero}
  \lean{AlgebraicGeometry.P1.gradeZeroAlgEquiv}
  \leanok
  The inclusion of constants \(k \to A_0\) into the degree-zero part of \(A\) is an
  isomorphism of \(k\)-algebras.
\end{definition}
\begin{proof}
  \leanok
  It is injective because a nonzero constant polynomial is nonzero, and surjective because a
  polynomial which is homogeneous of degree zero is constant.
\end{proof}

\begin{definition}[Structure morphism]
  \label{def:P1_structureMap}
  \lean{AlgebraicGeometry.P1.structureMap}
  \uses{def:P1, def:P1_gradeZero}
  \leanok
  The structure morphism \(\PP^1 \to \Spec k\) is the canonical morphism
  \(\Proj A \to \Spec A_0\) followed by the isomorphism \(\Spec A_0 \cong \Spec k\) induced by
  \ref{def:P1_gradeZero}.
\end{definition}

\begin{theorem}[Properness of the projective line]
  \label{thm:P1_proper}
  \lean{AlgebraicGeometry.P1.isProper_structureMap}
  \uses{def:P1_structureMap}
  \leanok
  \(\PP^1\) is proper over \(k\).
\end{theorem}
\begin{proof}
  \leanok
  \uses{def:P1_gradeZero}
  \(A\) is a finitely generated algebra over its degree-zero part \(A_0 = k\) (generated by
  \(X_0\) and \(X_1\)), and the Proj of a graded ring which is finitely generated over its
  degree-zero part is proper over the spectrum of the degree-zero part. Composing with the
  isomorphism \(\Spec A_0 \cong \Spec k\) preserves properness.
\end{proof}

\begin{definition}[The projective line as a scheme over \(k\)]
  \label{def:P1_asOver}
  \lean{AlgebraicGeometry.P1.asOver}
  \uses{def:P1_structureMap}
  \leanok
  \(\PP^1\), regarded as a scheme over \(k\) via the structure morphism of
  \ref{def:P1_structureMap}.
\end{definition}

\begin{definition}[The standard charts]
  \label{def:P1_chartOpen}
  \lean{AlgebraicGeometry.P1.chartOpen}
  \uses{def:P1}
  \leanok
  For \(i \in \{0,1\}\), the chart \(U_i = D_+(X_i) \subseteq \PP^1\) is the basic open locus
  where the homogeneous coordinate \(X_i\) does not vanish.
\end{definition}

\begin{lemma}[The charts are affine]
  \label{lem:P1_chartOpen_affine}
  \lean{AlgebraicGeometry.P1.isAffineOpen_chartOpen}
  \uses{def:P1_chartOpen}
  \leanok
  Each chart \(U_i\) is an affine open of \(\PP^1\).
\end{lemma}
\begin{proof}
  \leanok
  \(X_i\) is homogeneous of degree \(1 > 0\), and the basic open of \(\Proj\) at a homogeneous
  element of positive degree is affine.
\end{proof}

\begin{lemma}[The charts cover]
  \label{lem:P1_chartOpen_sup}
  \lean{AlgebraicGeometry.P1.chartOpen_sup}
  \uses{def:P1_chartOpen}
  \leanok
  \(U_0 \cup U_1 = \PP^1\).
\end{lemma}
\begin{proof}
  \leanok
  A point of \(\Proj A\) is a homogeneous prime ideal \(\mathfrak p\) of \(A\) not containing
  the irrelevant ideal \(A_+\). Every monomial of positive degree is divisible by \(X_0\) or
  by \(X_1\), so if both \(X_0 \in \mathfrak p\) and \(X_1 \in \mathfrak p\) then
  \(A_+ \subseteq \mathfrak p\), a contradiction. Hence every point avoids some \(X_i\), i.e.\
  lies in \(U_0 \cup U_1\).
\end{proof}

\begin{lemma}[The chart overlap]
  \label{lem:P1_chartOpen_inf}
  \lean{AlgebraicGeometry.P1.chartOpen_inf}
  \uses{def:P1_chartOpen}
  \leanok
  \(U_0 \cap U_1 = D_+(X_0 X_1)\).
\end{lemma}
\begin{proof}
  \leanok
  A homogeneous prime avoids the product \(X_0 X_1\) if and only if it avoids both factors.
\end{proof}

\begin{definition}[Affine models of the charts]
  \label{def:P1_chartIota}
  \lean{AlgebraicGeometry.P1.chartι}
  \uses{def:P1, def:P1_chartOpen}
  \leanok
  For \(i \in \{0,1\}\), the canonical morphism \(\Spec A_{(X_i)} \to \PP^1\) associated with
  the basic open at the degree-one element \(X_i\); it is an open immersion.
\end{definition}

\begin{lemma}[Image of the affine model]
  \label{lem:P1_chartIota_range}
  \lean{AlgebraicGeometry.P1.opensRange_chartι}
  \uses{def:P1_chartIota, def:P1_chartOpen}
  \leanok
  The image of \(\Spec A_{(X_i)} \to \PP^1\) is the chart \(U_i\); thus
  \(U_i \cong \Spec A_{(X_i)}\).
\end{lemma}
\begin{proof}
  \leanok
  This is the standard description of the basic opens of \(\Proj\): the homogeneous primes
  avoiding \(X_i\) correspond to the primes of the degree-zero localization \(A_{(X_i)}\).
\end{proof}

\begin{lemma}[The affine models lie over \(k\)]
  \label{lem:P1_chartIota_structureMap}
  \lean{AlgebraicGeometry.P1.chartι_structureMap}
  \uses{def:P1_chartIota, def:P1_structureMap}
  \leanok
  The composite \(\Spec A_{(X_i)} \to \PP^1 \to \Spec k\) is the morphism induced by the
  \(k\)-algebra structure of \(A_{(X_i)}\); that is, the affine models are morphisms of
  schemes over \(k\).
\end{lemma}
\begin{proof}
  \leanok
  \uses{def:P1_gradeZero}
  The canonical morphism \(\Proj A \to \Spec A_0\), restricted to the affine chart, is induced
  by the ring homomorphism \(A_0 \to A_{(X_i)}\); composing with the isomorphism
  \(k \cong A_0\) of \ref{def:P1_gradeZero} gives the claim.
\end{proof}

\begin{definition}[Chart coordinates]
  \label{def:P1_chartCoord}
  \lean{AlgebraicGeometry.P1.chartCoord}
  \uses{def:P1}
  \leanok
  For \(i, j \in \{0,1\}\), the chart coordinate is the degree-zero fraction
  \(t_{ij} = X_j / X_i \in A_{(X_i)}\).
\end{definition}

\begin{definition}[Dehomogenization]
  \label{def:P1_dehomogenize}
  \lean{AlgebraicGeometry.P1.dehomogenize}
  \leanok
  For \(i \in \{0,1\}\), the dehomogenization at \(X_i\) is the \(k\)-algebra homomorphism
  \(\delta_i : k[X_0, X_1] \to k[t]\) with \(\delta_i(X_i) = 1\) and \(\delta_i(X_j) = t\) for
  \(j \neq i\).
\end{definition}

\begin{definition}[From polynomials to the chart ring]
  \label{def:P1_polyToAway}
  \lean{AlgebraicGeometry.P1.polyToAway}
  \uses{def:P1_chartCoord}
  \leanok
  For \(i, j \in \{0,1\}\), the \(k\)-algebra homomorphism \(k[t] \to A_{(X_i)}\) given by
  evaluation at the chart coordinate, \(t \mapsto t_{ij}\).
\end{definition}

\begin{definition}[From the chart ring to polynomials]
  \label{def:P1_awayToPoly}
  \lean{AlgebraicGeometry.P1.awayToPoly}
  \uses{def:P1_dehomogenize}
  \leanok
  For \(i \in \{0,1\}\), the \(k\)-algebra homomorphism \(A_{(X_i)} \to k[t]\) sending a
  degree-zero fraction \(a / X_i^n\) (with \(a\) homogeneous of degree \(n\)) to
  \(\delta_i(a)\).
\end{definition}
\begin{proof}
  \leanok
  Well-definedness: \(\delta_i(X_i) = 1\) is a unit of \(k[t]\), so \(\delta_i\) extends to
  the localization \(A_{X_i} \to k[t]\) by the universal property of localization; on the
  degree-zero subring its value on \(a / X_i^n\) is
  \(\delta_i(a)\,\delta_i(X_i)^{-n} = \delta_i(a)\).
\end{proof}

\begin{lemma}[One composite is the identity]
  \label{lem:P1_awayToPoly_polyToAway}
  \lean{AlgebraicGeometry.P1.awayToPoly_comp_polyToAway}
  \uses{def:P1_polyToAway, def:P1_awayToPoly}
  \leanok
  For \(i \neq j\), the composite \(k[t] \to A_{(X_i)} \to k[t]\) of
  \ref{def:P1_polyToAway} followed by \ref{def:P1_awayToPoly} is the identity.
\end{lemma}
\begin{proof}
  \leanok
  Both sides are \(k\)-algebra homomorphisms on \(k[t]\), so it suffices to check the
  generator: \(t \mapsto t_{ij} = X_j / X_i \mapsto \delta_i(X_j) = t\).
\end{proof}

\begin{lemma}[Monomial fractions are powers of the coordinate]
  \label{lem:P1_away_monomial}
  \lean{AlgebraicGeometry.P1.away_mk_prod_eq}
  \uses{def:P1_chartCoord}
  \leanok
  Let \(i \neq j\) and let \(e_0, e_1 \ge 0\) with \(e_0 + e_1 = a\). Then in \(A_{(X_i)}\)
  \[
    \frac{X_0^{e_0} X_1^{e_1}}{X_i^{a}} \;=\; t_{ij}^{\,e_j}.
  \]
\end{lemma}
\begin{proof}
  \leanok
  Since \(X_0^{e_0} X_1^{e_1} = X_i^{e_i} X_j^{e_j}\) and \(a = e_i + e_j\), the left-hand
  side is \(X_i^{e_i} X_j^{e_j} / X_i^{e_i + e_j} = X_j^{e_j} / X_i^{e_j}\), which is
  \(t_{ij}^{\,e_j}\).
\end{proof}

\begin{lemma}[Surjectivity onto the chart ring]
  \label{lem:P1_polyToAway_surjective}
  \lean{AlgebraicGeometry.P1.polyToAway_surjective}
  \uses{def:P1_polyToAway}
  \leanok
  For \(i \neq j\), the homomorphism \(k[t] \to A_{(X_i)}\) of \ref{def:P1_polyToAway} is
  surjective.
\end{lemma}
\begin{proof}
  \leanok
  \uses{lem:P1_away_monomial, def:P1_gradeZero}
  Since \(A\) is generated over \(A_0\) by the degree-one elements \(X_0, X_1\), the
  degree-zero localization \(A_{(X_i)}\) is spanned, as a module over \(A_0 = k\)
  (\ref{def:P1_gradeZero}), by the fractions \(m / X_i^a\) where \(m\) ranges over the
  monomials \(X_0^{e_0} X_1^{e_1}\) of degree \(a\). By \ref{lem:P1_away_monomial} each such
  fraction equals \(t_{ij}^{\,e_j}\), which is the image of \(t^{e_j}\). The image of a ring
  homomorphism is closed under addition and under multiplication by scalars from \(k\), so it
  is all of \(A_{(X_i)}\).
\end{proof}

\begin{lemma}[The other composite is the identity]
  \label{lem:P1_polyToAway_awayToPoly}
  \lean{AlgebraicGeometry.P1.polyToAway_comp_awayToPoly}
  \uses{def:P1_polyToAway, def:P1_awayToPoly}
  \leanok
  For \(i \neq j\), the composite \(A_{(X_i)} \to k[t] \to A_{(X_i)}\) of
  \ref{def:P1_awayToPoly} followed by \ref{def:P1_polyToAway} is the identity.
\end{lemma}
\begin{proof}
  \leanok
  \uses{lem:P1_polyToAway_surjective, lem:P1_awayToPoly_polyToAway}
  By \ref{lem:P1_polyToAway_surjective} every element of \(A_{(X_i)}\) is of the form
  \(p(t_{ij})\) with \(p \in k[t]\), and on such elements the claim follows from
  \ref{lem:P1_awayToPoly_polyToAway}.
\end{proof}

\begin{definition}[The chart identification]
  \label{def:P1_awayAlgEquiv}
  \lean{AlgebraicGeometry.P1.awayAlgEquiv}
  \uses{def:P1_polyToAway, def:P1_awayToPoly}
  \leanok
  For \(i \neq j\), the isomorphism of \(k\)-algebras
  \(\alpha_i : A_{(X_i)} \cong k[t]\) given by the mutually inverse homomorphisms of
  \ref{def:P1_awayToPoly} and \ref{def:P1_polyToAway}; the chart coordinate \(t_{ij}\)
  corresponds to \(t\).
\end{definition}
\begin{proof}
  \leanok
  \uses{lem:P1_awayToPoly_polyToAway, lem:P1_polyToAway_awayToPoly}
  The two composites are the identity by \ref{lem:P1_awayToPoly_polyToAway} and
  \ref{lem:P1_polyToAway_awayToPoly}.
\end{proof}

\section{The chart overlap and Laurent polynomials}

Write \(k[T, T^{-1}]\) for the ring of Laurent polynomials over \(k\).

\begin{lemma}[Laurent span]
  \label{lem:laurent_span}
  \lean{LaurentPolynomial.exists_toLaurent_add_aeval}
  \leanok
  Every Laurent polynomial \(f \in R[T, T^{-1}]\) over a commutative (semi)ring \(R\) can be
  written as \(f = p(T) + q(T^{-1})\) with \(p, q \in R[t]\) ordinary polynomials.
\end{lemma}
\begin{proof}
  \leanok
  Split \(f\) into its monomials and collect those with nonnegative exponent into \(p\) and
  those with negative exponent into \(q\).
\end{proof}

\begin{definition}[Restriction from the left chart]
  \label{def:P1_awayToOverlapLeft}
  \lean{AlgebraicGeometry.P1.awayToOverlapLeft}
  \uses{def:P1}
  \leanok
  The canonical ring homomorphism \(\lambda : A_{(X_0)} \to A_{(X_0 X_1)}\), sending
  \(a / X_0^{\,n}\) to \(a X_1^{\,n} / (X_0 X_1)^{n}\).
\end{definition}

\begin{definition}[Restriction from the right chart]
  \label{def:P1_awayToOverlapRight}
  \lean{AlgebraicGeometry.P1.awayToOverlapRight}
  \uses{def:P1}
  \leanok
  The canonical ring homomorphism \(\rho : A_{(X_1)} \to A_{(X_0 X_1)}\), sending
  \(a / X_1^{\,n}\) to \(a X_0^{\,n} / (X_0 X_1)^{n}\).
\end{definition}

\begin{lemma}[The two coordinates are mutually inverse on the overlap]
  \label{lem:P1_coord_mul}
  \lean{AlgebraicGeometry.P1.awayToOverlap_mul_eq_one}
  \uses{def:P1_awayToOverlapLeft, def:P1_awayToOverlapRight, def:P1_chartCoord}
  \leanok
  In \(A_{(X_0 X_1)}\) one has \(\lambda(t_{01}) \cdot \rho(t_{10}) = 1\).
\end{lemma}
\begin{proof}
  \leanok
  Computing with the formulas of \ref{def:P1_awayToOverlapLeft} and
  \ref{def:P1_awayToOverlapRight}:
  \(\lambda(t_{01}) = X_1^2 / (X_0 X_1)\) and \(\rho(t_{10}) = X_0^2 / (X_0 X_1)\), so
  \[
    \lambda(t_{01}) \cdot \rho(t_{10})
      = \frac{X_0^2 X_1^2}{(X_0 X_1)^2} = 1 .
  \]
\end{proof}

\begin{lemma}[The overlap ring is a localization of the chart ring]
  \label{lem:P1_overlap_isLocalization}
  \lean{AlgebraicGeometry.P1.instAwayAwayNatMvPolynomialFinOfNatSubmoduleHomogeneousSubmoduleXChartCoordHMul}
  \uses{def:P1_awayToOverlapLeft, def:P1_chartCoord}
  \leanok
  The homomorphism \(\lambda\) of \ref{def:P1_awayToOverlapLeft} identifies
  \(A_{(X_0 X_1)}\) with the localization of \(A_{(X_0)}\) at the powers of the chart
  coordinate \(t_{01}\).
\end{lemma}
\begin{proof}
  \leanok
  This is a general property of degree-zero localizations: for homogeneous elements \(f, g\)
  of degree one, \(A_{(fg)}\) is the localization of \(A_{(f)}\) away from the degree-zero
  fraction \(g/f\). Here \(f = X_0\), \(g = X_1\) and \(g/f = t_{01}\).
\end{proof}

\begin{definition}[The overlap identification]
  \label{def:P1_overlapAlgEquiv}
  \lean{AlgebraicGeometry.P1.overlapAlgEquiv}
  \uses{def:P1}
  \leanok
  The isomorphism of \(k\)-algebras \(\theta : A_{(X_0 X_1)} \cong k[T, T^{-1}]\), with
  \(T\) the image of the coordinate \(t_{01}\).
\end{definition}
\begin{proof}
  \leanok
  \uses{lem:P1_overlap_isLocalization, def:P1_awayAlgEquiv}
  By \ref{lem:P1_overlap_isLocalization}, \(A_{(X_0X_1)}\) is the localization of
  \(A_{(X_0)}\) at the powers of \(t_{01}\). The isomorphism
  \(\alpha_0 : A_{(X_0)} \cong k[t]\) of \ref{def:P1_awayAlgEquiv} carries \(t_{01}\) to
  \(t\), hence induces an isomorphism of the two localizations
  \(A_{(X_0X_1)} \cong k[t]_{t} = k[T, T^{-1}]\) (localizations at corresponding
  multiplicative sets are canonically isomorphic). The map is \(k\)-linear because
  \(\alpha_0\) is.
\end{proof}

\begin{lemma}[Left restriction is \(t \mapsto T\)]
  \label{lem:P1_overlap_res_left}
  \lean{AlgebraicGeometry.P1.overlapAlgEquiv_awayToOverlapLeft}
  \uses{def:P1_overlapAlgEquiv, def:P1_awayAlgEquiv, def:P1_awayToOverlapLeft}
  \leanok
  For every \(z \in A_{(X_0)}\) one has \(\theta(\lambda(z)) = p(T)\), where
  \(p = \alpha_0(z) \in k[t]\). In other words, under the identifications
  \(\alpha_0\) and \(\theta\), the restriction \(\lambda\) is the canonical inclusion
  \(k[t] \to k[T, T^{-1}]\), \(t \mapsto T\).
\end{lemma}
\begin{proof}
  \leanok
  \uses{lem:P1_overlap_isLocalization}
  By construction (\ref{def:P1_overlapAlgEquiv}), \(\theta\) is the comparison of
  localizations induced by \(\alpha_0\); under such a comparison the localization map
  \(\lambda\) of \ref{lem:P1_overlap_isLocalization} corresponds to the localization map
  \(k[t] \to k[T, T^{-1}]\).
\end{proof}

\begin{lemma}[Right restriction is \(t \mapsto T^{-1}\)]
  \label{lem:P1_overlap_res_right}
  \lean{AlgebraicGeometry.P1.overlapAlgEquiv_awayToOverlapRight}
  \uses{def:P1_overlapAlgEquiv, def:P1_awayAlgEquiv, def:P1_awayToOverlapRight}
  \leanok
  For every \(z \in A_{(X_1)}\) one has \(\theta(\rho(z)) = q(T^{-1})\), where
  \(q = \alpha_1(z) \in k[t]\). In other words, under the identifications \(\alpha_1\) and
  \(\theta\), the restriction \(\rho\) is the homomorphism \(k[t] \to k[T, T^{-1}]\),
  \(t \mapsto T^{-1}\).
\end{lemma}
\begin{proof}
  \leanok
  \uses{lem:P1_coord_mul, lem:P1_overlap_res_left, lem:P1_polyToAway_surjective}
  By \ref{lem:P1_overlap_res_left} the image of \(\lambda(t_{01})\) under \(\theta\) is
  \(T\), and by \ref{lem:P1_coord_mul} \(\lambda(t_{01})\,\rho(t_{10}) = 1\); applying
  \(\theta\) gives \(\theta(\rho(t_{10})) = T^{-1}\). By
  \ref{lem:P1_polyToAway_surjective} (for the indices \((1,0)\)) every \(z \in A_{(X_1)}\)
  is a polynomial in \(t_{10}\) with coefficients in \(k\); both sides of the claim are
  \(k\)-algebra homomorphisms in \(z\), and they agree on \(t_{10}\).
\end{proof}

\begin{lemma}[Laurent span on the overlap ring]
  \label{lem:P1_overlap_span}
  \lean{AlgebraicGeometry.P1.exists_awayToOverlap_add}
  \uses{def:P1_awayToOverlapLeft, def:P1_awayToOverlapRight}
  \leanok
  Every \(w \in A_{(X_0 X_1)}\) can be written as \(w = \lambda(z_0) + \rho(z_1)\) with
  \(z_0 \in A_{(X_0)}\) and \(z_1 \in A_{(X_1)}\).
\end{lemma}
\begin{proof}
  \leanok
  \uses{lem:laurent_span, def:P1_overlapAlgEquiv, lem:P1_overlap_res_left,
    lem:P1_overlap_res_right, def:P1_awayAlgEquiv}
  Write \(\theta(w) = p(T) + q(T^{-1})\) by \ref{lem:laurent_span} and set
  \(z_0 = \alpha_0^{-1}(p)\), \(z_1 = \alpha_1^{-1}(q)\). By
  \ref{lem:P1_overlap_res_left} and \ref{lem:P1_overlap_res_right},
  \(\theta(\lambda(z_0) + \rho(z_1)) = p(T) + q(T^{-1}) = \theta(w)\), and \(\theta\) is
  injective.
\end{proof}

\section{Section rings and the Laurent chart pair}

\begin{lemma}[The overlap is affine]
  \label{lem:P1_overlap_affine}
  \lean{AlgebraicGeometry.P1.isAffineOpen_overlap}
  \uses{lem:P1_chartOpen_inf}
  \leanok
  The chart overlap \(U_0 \cap U_1 = D_+(X_0 X_1)\) is an affine open of \(\PP^1\).
\end{lemma}
\begin{proof}
  \leanok
  \(X_0 X_1\) is homogeneous of degree \(2 > 0\), and the basic open of \(\Proj\) at a
  homogeneous element of positive degree is affine.
\end{proof}

\begin{definition}[Sections on the left chart]
  \label{def:P1_chartSectionsEquiv0}
  \lean{AlgebraicGeometry.P1.chartSectionsEquiv₀}
  \uses{def:P1_chartOpen, def:P1_chartIota, def:P1_awayAlgEquiv}
  \leanok
  The ring isomorphism \(\gamma_0 : \Gamma(\PP^1, U_0) \cong k[t]\): the canonical
  identification \(\Gamma(\PP^1, D_+(X_0)) \cong A_{(X_0)}\) (sections of \(\Proj\) over a
  basic open; equivalently, \(U_0 = \Spec A_{(X_0)}\) via \ref{def:P1_chartIota}) followed by
  \(\alpha_0\) of \ref{def:P1_awayAlgEquiv}.
\end{definition}

\begin{definition}[Sections on the right chart]
  \label{def:P1_chartSectionsEquiv1}
  \lean{AlgebraicGeometry.P1.chartSectionsEquiv₁}
  \uses{def:P1_chartOpen, def:P1_chartIota, def:P1_awayAlgEquiv}
  \leanok
  The ring isomorphism \(\gamma_1 : \Gamma(\PP^1, U_1) \cong k[t]\), defined as in
  \ref{def:P1_chartSectionsEquiv0} with the identification
  \(\Gamma(\PP^1, D_+(X_1)) \cong A_{(X_1)}\) followed by \(\alpha_1\).
\end{definition}

\begin{definition}[Sections on the overlap]
  \label{def:P1_overlapSectionsEquiv}
  \lean{AlgebraicGeometry.P1.overlapSectionsEquiv}
  \uses{lem:P1_chartOpen_inf, def:P1_overlapAlgEquiv}
  \leanok
  The ring isomorphism \(\gamma_{01} : \Gamma(\PP^1, U_0 \cap U_1) \cong k[T, T^{-1}]\): the
  canonical identification \(\Gamma(\PP^1, D_+(X_0X_1)) \cong A_{(X_0X_1)}\) followed by
  \(\theta\) of \ref{def:P1_overlapAlgEquiv}.
\end{definition}

\begin{lemma}[Restriction of sections from the left chart]
  \label{lem:P1_res_left}
  \lean{AlgebraicGeometry.P1.overlapSectionsEquiv_res_left}
  \uses{def:P1_chartSectionsEquiv0, def:P1_overlapSectionsEquiv}
  \leanok
  For every section \(a \in \Gamma(\PP^1, U_0)\),
  \[
    \gamma_{01}\bigl(a|_{U_0 \cap U_1}\bigr) = p(T), \qquad p = \gamma_0(a);
  \]
  under the identifications, restriction from the left chart to the overlap is
  \(t \mapsto T\).
\end{lemma}
\begin{proof}
  \leanok
  \uses{lem:P1_overlap_res_left}
  The identifications \(\Gamma(\PP^1, D_+(f)) \cong A_{(f)}\) are natural with respect to
  inclusions of basic opens: restriction of sections from \(D_+(X_0)\) to \(D_+(X_0X_1)\)
  corresponds to the ring homomorphism \(\lambda\). Conclude by
  \ref{lem:P1_overlap_res_left}.
\end{proof}

\begin{lemma}[Restriction of sections from the right chart]
  \label{lem:P1_res_right}
  \lean{AlgebraicGeometry.P1.overlapSectionsEquiv_res_right}
  \uses{def:P1_chartSectionsEquiv1, def:P1_overlapSectionsEquiv}
  \leanok
  For every section \(b \in \Gamma(\PP^1, U_1)\),
  \[
    \gamma_{01}\bigl(b|_{U_0 \cap U_1}\bigr) = q(T^{-1}), \qquad q = \gamma_1(b);
  \]
  under the identifications, restriction from the right chart to the overlap is
  \(t \mapsto T^{-1}\).
\end{lemma}
\begin{proof}
  \leanok
  \uses{lem:P1_overlap_res_right}
  As in \ref{lem:P1_res_left}, restriction of sections corresponds to \(\rho\); conclude by
  \ref{lem:P1_overlap_res_right}.
\end{proof}

\begin{definition}[Laurent chart pair]
  \label{def:LaurentChartPair}
  \lean{AlgebraicGeometry.LaurentChartPair}
  \uses{def:P1}
  \leanok
  A \emph{Laurent chart pair} on \(\PP^1\) consists of open subsets
  \(V_0, V_1, V_{01} \subseteq \PP^1\) with \(V_{01} \subseteq V_0\),
  \(V_{01} \subseteq V_1\) and \(V_0 \cap V_1 = V_{01}\), all three affine, covering the
  projective line (\(V_0 \cup V_1 = \PP^1\)), together with ring isomorphisms
  \[
    \Gamma_0 : \Gamma(\PP^1, V_0) \cong k[t], \quad
    \Gamma_1 : \Gamma(\PP^1, V_1) \cong k[t], \quad
    \Gamma_{01} : \Gamma(\PP^1, V_{01}) \cong k[T, T^{-1}],
  \]
  such that, under these identifications, restriction \(V_{01} \subseteq V_0\) is
  \(t \mapsto T\) and restriction \(V_{01} \subseteq V_1\) is \(t \mapsto T^{-1}\).
\end{definition}

\begin{theorem}[Laurent span for a chart pair]
  \label{thm:laurentChartPair_span}
  \lean{AlgebraicGeometry.LaurentChartPair.exists_res_add_res}
  \uses{def:LaurentChartPair}
  \leanok
  Let \((V_0, V_1, V_{01}, \Gamma_0, \Gamma_1, \Gamma_{01})\) be a Laurent chart pair on
  \(\PP^1\). Every section \(z \in \Gamma(\PP^1, V_{01})\) is the sum of a restriction from
  \(V_0\) and a restriction from \(V_1\): there are \(a \in \Gamma(\PP^1, V_0)\) and
  \(b \in \Gamma(\PP^1, V_1)\) with \(z = a|_{V_{01}} + b|_{V_{01}}\).
\end{theorem}
\begin{proof}
  \leanok
  \uses{lem:laurent_span}
  Write \(\Gamma_{01}(z) = p(T) + q(T^{-1})\) by \ref{lem:laurent_span} and set
  \(a = \Gamma_0^{-1}(p)\), \(b = \Gamma_1^{-1}(q)\). By the two restriction conditions,
  \(\Gamma_{01}(a|_{V_{01}} + b|_{V_{01}}) = p(T) + q(T^{-1}) = \Gamma_{01}(z)\), and
  \(\Gamma_{01}\) is injective.
\end{proof}

\begin{definition}[The canonical Laurent chart pair]
  \label{def:P1_laurentChartPair}
  \lean{AlgebraicGeometry.P1.laurentChartPair}
  \uses{def:LaurentChartPair, def:P1_chartOpen, def:P1_chartSectionsEquiv0,
    def:P1_chartSectionsEquiv1, def:P1_overlapSectionsEquiv}
  \leanok
  The Laurent chart pair on \(\PP^1\) given by the standard charts: \(V_0 = U_0\),
  \(V_1 = U_1\), \(V_{01} = D_+(X_0X_1)\), with the identifications \(\gamma_0, \gamma_1,
  \gamma_{01}\) of \ref{def:P1_chartSectionsEquiv0}--\ref{def:P1_overlapSectionsEquiv}.
\end{definition}
\begin{proof}
  \leanok
  \uses{lem:P1_chartOpen_affine, lem:P1_chartOpen_sup, lem:P1_chartOpen_inf,
    lem:P1_overlap_affine, lem:P1_res_left, lem:P1_res_right}
  The overlap conditions and affineness are \ref{lem:P1_chartOpen_inf},
  \ref{lem:P1_chartOpen_affine} and \ref{lem:P1_overlap_affine}; the covering condition is
  \ref{lem:P1_chartOpen_sup}; the two restriction conditions are \ref{lem:P1_res_left} and
  \ref{lem:P1_res_right}.
\end{proof}

\section{A finite morphism to the projective line}

\begin{theorem}[Zariski's main theorem, cancellation form]
  \label{thm:zmt_cancellation}
  \lean{AlgebraicGeometry.IsFinite.of_locallyQuasiFinite_of_comp}
  \leanok
  Let \(\pi : X \to Y\) and \(g : Y \to S\) be morphisms of schemes with \(\pi\) locally
  quasi-finite, \(g\) separated and \(g \circ \pi\) proper. Then \(\pi\) is finite.
\end{theorem}
\begin{proof}
  \leanok
  By the cancellation property of proper morphisms (using that \(g\) is separated), \(\pi\)
  is proper. A proper, locally quasi-finite morphism is finite (Zariski's main theorem).
\end{proof}

\begin{lemma}[Locally quasi-finite maps to the projective line are finite]
  \label{lem:finite_toP1_of_lqf}
  \lean{AlgebraicGeometry.isFinite_toP1_of_locallyQuasiFinite}
  \uses{def:P1_structureMap}
  \leanok
  Let \(X\) be a scheme, proper over \(k\), and let \(\pi : X \to \PP^1\) be a locally
  quasi-finite morphism over \(k\) (i.e.\ its composite with the structure morphism of
  \(\PP^1\) is the structure morphism of \(X\)). Then \(\pi\) is finite.
\end{lemma}
\begin{proof}
  \leanok
  \uses{thm:zmt_cancellation, thm:P1_proper}
  The structure morphism of \(\PP^1\) is proper (\ref{thm:P1_proper}), in particular
  separated, and the composite of \(\pi\) with it is the structure morphism of \(X\), which
  is proper. Apply \ref{thm:zmt_cancellation}.
\end{proof}

\begin{lemma}[Reduction to local quasi-finiteness]
  \label{lem:exists_finite_of_lqf}
  \lean{AlgebraicGeometry.exists_isFinite_toP1_of_locallyQuasiFinite}
  \uses{def:P1_structureMap}
  \leanok
  Let \(C\) be a scheme, proper over \(k\). If there exists a locally quasi-finite morphism
  \(C \to \PP^1\) over \(k\), then there exists a finite morphism \(C \to \PP^1\) over \(k\).
\end{lemma}
\begin{proof}
  \leanok
  \uses{lem:finite_toP1_of_lqf}
  Immediate from \ref{lem:finite_toP1_of_lqf}: the given locally quasi-finite morphism is
  itself finite.
\end{proof}

\begin{theorem}[Existence of a finite morphism to the projective line]
  \label{thm:exists_finite_toP1}
  \lean{AlgebraicGeometry.exists_isFinite_toP1}
  \uses{def:P1, def:P1_structureMap}
  \leanok
  Let \(C\) be a smooth curve over \(k\) (proper, smooth of relative dimension one,
  geometrically irreducible). Then there exists a finite morphism \(\pi : C \to \PP^1\) over
  \(k\).
\end{theorem}

By \ref{lem:exists_finite_of_lqf}, it suffices to produce a locally quasi-finite morphism
\(C \to \PP^1\) over \(k\); such a morphism arises from a rational function on \(C\)
transcendental over \(k\), extended to all of \(C\) using that the local rings of \(C\) at
closed points are discrete valuation rings.

\begin{lemma}[Preimages of the charts are affine]
  \label{lem:preimage_chart_affine}
  \lean{AlgebraicGeometry.isAffineOpen_preimage_chartOpen}
  \uses{def:P1_chartOpen}
  \leanok
  Let \(\pi : Y \to \PP^1\) be an affine morphism (for instance a finite one). Then
  \(\pi^{-1}(U_i)\) is an affine open of \(Y\) for \(i \in \{0,1\}\).
\end{lemma}
\begin{proof}
  \leanok
  \uses{lem:P1_chartOpen_affine}
  The charts are affine (\ref{lem:P1_chartOpen_affine}), and the preimage of an affine open
  under an affine morphism is affine.
\end{proof}

\begin{lemma}[Preimages of the charts cover]
  \label{lem:preimage_chart_sup}
  \lean{AlgebraicGeometry.preimage_chartOpen_sup}
  \uses{def:P1_chartOpen}
  \leanok
  For any morphism \(\pi : Y \to \PP^1\), the preimages of the two standard charts cover
  \(Y\): \(\pi^{-1}(U_0) \cup \pi^{-1}(U_1) = Y\).
\end{lemma}
\begin{proof}
  \leanok
  \uses{lem:P1_chartOpen_sup}
  Preimages of a cover form a cover (\ref{lem:P1_chartOpen_sup}).
\end{proof}

\begin{lemma}[Module-finiteness over the charts]
  \label{lem:finite_app_chart}
  \lean{AlgebraicGeometry.finite_app_chartOpen}
  \uses{def:P1_chartOpen}
  \leanok
  Let \(\pi : Y \to \PP^1\) be a finite morphism. For \(i \in \{0,1\}\), the induced ring
  homomorphism \(\Gamma(\PP^1, U_i) \to \Gamma(Y, \pi^{-1}(U_i))\) is module-finite. (By
  \ref{def:P1_chartSectionsEquiv0}--\ref{def:P1_chartSectionsEquiv1}, the source is the
  polynomial ring \(k[t]\).)
\end{lemma}
\begin{proof}
  \leanok
  \uses{lem:P1_chartOpen_affine}
  A finite morphism induces module-finite maps on sections over affine opens; the charts are
  affine by \ref{lem:P1_chartOpen_affine}.
\end{proof}

\begin{lemma}[Module-finiteness over the overlap]
  \label{lem:finite_app_overlap}
  \lean{AlgebraicGeometry.finite_app_overlap}
  \uses{lem:P1_chartOpen_inf}
  \leanok
  Let \(\pi : Y \to \PP^1\) be a finite morphism. The induced ring homomorphism
  \(\Gamma(\PP^1, U_0 \cap U_1) \to \Gamma(Y, \pi^{-1}(U_0 \cap U_1))\) is module-finite.
  (By \ref{def:P1_overlapSectionsEquiv}, the source is the Laurent polynomial ring
  \(k[T,T^{-1}]\).)
\end{lemma}
\begin{proof}
  \leanok
  \uses{lem:P1_overlap_affine}
  As in \ref{lem:finite_app_chart}, with the overlap affine by \ref{lem:P1_overlap_affine}.
\end{proof}
